\section{Additional delayed-observation calibration panel}
\label{app:t1_validation}
This supplementary CPU study evaluates the operational pending-outcome adapter
(ADAPTER), a completed-prefix construction (CPREFIX), and a naive completed-only
construction (NAIVE). It uses 28,000 programs, each containing four independently
generated trial paths, for 112,000 trials at an enrollment horizon of 2,000 pairs.
Constructions share the same path within a trial; cells use separate random streams.
The primary observation schedule processes events in tick batches, including the
fixed finalization window. This is a synthetic calibration study, not a live-agent
experiment. Namespace-0 coordinates had prior development exposure and are reported
as replay evidence, not a fresh confirmatory holdout.

The eight cells cross four outcome laws with non-informative (odd cell numbers)
and informative (even cell numbers) delay. Their hierarchy/success means are
$(0,0)$ for C1--C2, $(0,0.25)$ for C3--C4, $(0.40,-0.03)$ for C5--C6,
and $(0.45,0.20)$ for C7--C8. Deployment requires a positive hierarchy contrast
and success noninferiority with margin $0.03$; C5--C6 are exactly at that
noninferiority boundary. The frozen allocations are 2,000 programs per cell
for C1, C2, C7 and C8, and 5,000 for C3--C6.

A false deployment is a deployment when the conjunction of hierarchy superiority
and success noninferiority is false. A false-harm decision retains the incumbent
when the hierarchy mean is nonnegative. Trial error is the union of those events;
program-family error is their union across its four trials. The nominal per-band,
per-trial union and per-program levels are respectively $0.00625$, $0.0125$ and
$0.05$. Ever-miscoverage and decision error are distinct readouts.

\begin{table}[ht]
\centering\small
\begin{tabular}{lrrrrrr}
\toprule
& Trials & \multicolumn{2}{c}{Trial errors} & \multicolumn{2}{c}{Hierarchy miscoverage} & Correct deploys\\
Cell & per cell & ADAPTER & NAIVE & ADAPTER & NAIVE & ADAPTER\\
\midrule
C1 & 8,000 & 0 & 2 & 0 & 2 & 0\\
C2 & 8,000 & 0 & 404 & 1 & 7,910 & 0\\
C3 & 20,000 & 0 & 23 & 0 & 23 & 0\\
C4 & 20,000 & 5 & 19,998 & 7 & 19,998 & 0\\
C5 & 20,000 & 0 & 0 & 0 & 0 & 0\\
C6 & 20,000 & 0 & 0 & 0 & 16,272 & 0\\
C7 & 8,000 & 0 & 0 & 0 & 2 & 8,000\\
C8 & 8,000 & 0 & 0 & 0 & 5,550 & 8,000\\
\bottomrule
\end{tabular}
\caption{Counts on committed trial paths. Each construction has the listed trial
denominator; constructions are paired, not pooled. Hierarchy miscoverage is an
indicator of any two-sided exclusion on the recorded path. These selected laws
supply finite-design calibration evidence, not a universal validity guarantee.}
\label{tab:t1_validation}
\end{table}

Under informative delay, NAIVE's program-family error is $376/2{,}000=0.188$
in C2 (marginal Wilson 95\% interval $[0.1715,0.2057]$) and $5{,}000/5{,}000$
in C4 ($[0.9992,1]$). ADAPTER's only trial errors are five false deployments in
C4: $5/20{,}000=0.00025$, with interval $[0.0001068,0.0005852]$.
All five records also flag eligible-look hierarchy miscoverage; five distinct
programs are affected. This nonzero count does not demonstrate violation of the
nominal error bound. CPREFIX has 15 false-harm records (C1: 1; C3: 10; C4: 4),
and NAIVE has 14 (C1: 2; C3: 12); all 29 retention records are included in the
corrected summaries. Intervals describe Monte Carlo uncertainty under the stated
generator, are marginal rather than simultaneous, and do not acquire anytime
coverage from the underlying monitoring rule.

ADAPTER abstains in every C1, C2, C3, C5 and C6 trial and in 19,995 of 20,000
C4 trials, while deploying correctly in every C7 and C8 trial. Thus the result
combines low observed errors with substantial abstention outside strong
alternatives. It does not identify a power curve or minimum detectable effect.
Zero deployments at the success boundary describe the full guarded procedure;
no hierarchy-only ablation is used to attribute that result solely to the guardrail.
C6 and C8 further show why low decision error cannot by itself establish valid
bands: NAIVE has extensive hierarchy miscoverage in both, despite no wrong
first decision in those cells. The complete-information H/D reference records
were checked for identities and hashes, but no reference-width or reference-power
comparison is claimed here.

\paragraph{Reproducibility and limitations.}
The accepted snapshot is \texttt{35ab9d1}, with independent aggregation and
provenance reports at \texttt{97ad956} in the repository. All 56 shards, 112 raw
file hashes, trial/program coordinates, and source/configuration pins were
reconciled; key event counts and Wilson intervals were independently recomputed
from committed records. This reproduction checks exported outcomes and flags,
not an independent regeneration of every latent path or band. The code archive
includes unchanged primary records and a standard-library table reproducer;
full original receipts and reference outputs remain linked to the exact Git snapshot.
The parent failed after the child completed; finalization was recovered without
repeating the scientific trials. Child supervision recorded 2,395.19 seconds and
a sampled peak tree RSS of 75.77 MiB. Continuous monitoring of the recovery stage
was not established. Original worker native/bytecode hashes were not re-executed
locally. The owner's later clarification at \texttt{03dc070} identifies the
discarded passes as earlier measurement-mode executions, consistent with the
existing development ledger. Their deleted records and unknown usage remain
unrecoverable; the disclosed exposure to decision counts is not undone. This
clarification does not establish the absence of unrecorded attempts.
The accepted completed-panel records remain unchanged. These limitations, prior
coordinate exposure and the finite simulated design restrict the claims above.
